\section{Introducción: por qué en 2025 se volvió a «tallar la Attention»}

Esta sección establece primero el contexto técnico de todo el episodio. Yang Songlin (杨松琳) habla de uno de los módulos más importantes de la arquitectura básica de los modelos grandes: la Attention. En los últimos años, la parte FFN del Transformer ya se modificó a fondo con MoE (Mixture of Experts, mezcla de expertos), y modelos como DeepSeek hicieron de MoE un consenso. La siguiente pieza que podría «tallarse» es la Attention. El motivo es directo: el contexto largo, el CoT largo, la Agentic AI y el decoding de textos largos pusieron en primer plano la carga de cómputo y de KV cache de la Full Attention.

Este episodio busca responder cuatro preguntas. Primera: ¿en qué consiste exactamente lo «linear» de la Linear Attention? Segunda: ¿por qué es importante el módulo KDA de Kimi Linear? Tercera: ¿qué compromisos implican, respectivamente, la atención lineal híbrida de Kimi, la atención dispersa de DeepSeek y el regreso de MiniMax M2 a la Full Attention? Cuarta: ¿por qué la investigación en arquitecturas no puede limitarse a la elegancia matemática y debe coordinarse con la afinidad con el hardware, los algoritmos paralelos, la multiplicación de matrices y la optimización de kernels?

\begin{importantbox}{Tesis central del episodio}
Cuando el crecimiento de los datos se desacelera, aumenta la demanda de contexto largo y el costo de inferencia se vuelve más pesado, la innovación en arquitecturas vuelve a ser decisiva. Las nuevas líneas de Attention no buscan sustituir el nombre Transformer, sino obtener, con los mismos FLOPs, una loss más baja, menos carga de KV cache y un decoding más eficiente en secuencias largas.
\end{importantbox}

\begin{knowledgebox}{Nota sobre la estrategia visual}
Este video es una entrevista con encuadre fijo, sin slides ni proyección de artículos. El texto no incluye fotogramas repetidos de los participantes; todas las figuras son diagramas conceptuales elaborados para esta nota, que explican la complejidad de la atención, KDA, las tres líneas de Attention, la Scaling Ladder, la arqueología de algoritmos y la afinidad con el hardware.
\end{knowledgebox}

\subsection{Resumen de la sección}

EP119 es una revisión panorámica de arquitecturas, no una entrevista común. Reúne en un mismo mapa las distintas apuestas que las empresas de modelos hicieron en 2025 sobre la Attention: Kimi apuesta por Linear/Hybrid, DeepSeek por Sparse y MiniMax M2 regresa por ahora a la Full Attention.
